\documentclass[11pt]{article}
% Source-PDF: 动态规划 DYNAMIC PROGRAMMING/动态规划加速原理之四边形不等式 - 赵爽.pdf
\usepackage{oipl-vn}

\title{Nguyên lý tăng tốc quy hoạch động bằng bất đẳng thức tứ giác}
\author{Zhao Shuang}
\date{}

\begin{document}
\maketitle

\origtitle{Principle of DP Acceleration by the Quadrangle Inequality}
\sourcepath{Dynamic Programming / Quadrangle Inequality Acceleration - Zhao Shuang.pdf}

\section{Lý thuyết cơ bản về bất đẳng thức tứ giác}

Trong các bài toán quy hoạch động, ta thường gặp một phương trình chuyển
trạng thái dạng sau:
\[
  m(i,j)=
  \begin{cases}
    \displaystyle
    \min_{i<k\le j}
      \{m(i,k-1)+m(k,j)+w(i,j)\}, & i<j,\\[0.8em]
    0, & i=j,\\
    \infty, & i>j.
  \end{cases}
  \tag{1.1}
\]
Chẳng hạn, bài toán ``cây con mẹ có chi phí nhỏ nhất'' cũng dùng công thức
này. Với các bài toán khác nhau, giá trị biên của $m(i,j)$ có thể khác nhau;
điều đó không ảnh hưởng đến phần thảo luận dưới đây.

Giả sử với $i\le i'<j\le j'$, ta có
\[
  w(i',j)\le w(i,j').
\]
Khi đó ta nói hàm $w$ có tính đơn điệu theo quan hệ bao hàm của đoạn. Ngoài
ra, nếu
\[
  w(i,j)+w(i',j')\le w(i',j)+w(i,j'),
  \tag{1.2}
\]
ta nói hàm $w$ thỏa mãn bất đẳng thức tứ giác.

Ví dụ, trong bài toán ``cây con mẹ có chi phí nhỏ nhất'', ta có
\[
  w(i,j)+w(i',j')=w(i',j)+w(i,j').
\]
Vì vậy hàm $w$ của bài toán đó thỏa mãn bất đẳng thức tứ giác. Có thể hình
dung bất đẳng thức này như sau: trong một tứ giác, tổng trọng số tại hai đầu
của một đường chéo không lớn hơn tổng trọng số tại hai đầu của đường chéo còn
lại.

Từ đây về sau, ta xét bài toán trong điều kiện $w$ vừa có tính đơn điệu theo
quan hệ bao hàm của đoạn, vừa thỏa mãn bất đẳng thức tứ giác.

\paragraph{Định lý 1.}
Nếu hàm $w$ thỏa mãn các điều kiện trên, thì hàm $m$ cũng thỏa mãn bất đẳng
thức tứ giác, tức là
\[
  m(i,j)+m(i',j')\le m(i',j)+m(i,j'),
  \qquad i\le i'<j\le j'.
\]

\paragraph{Chứng minh.}
Khi $i=i'$ hoặc $j=j'$, hiển nhiên có
\[
  m(i,j)+m(i',j')=m(i',j)+m(i,j'),
\]
nên bất đẳng thức đúng.

Với các trường hợp còn lại, ta chứng minh quy nạp theo
$l=j'-i$. Có hai tình huống cần xét.

\paragraph{Trường hợp 1: $i<i'=j<j'$.}
Khi đó bất đẳng thức tứ giác trở thành dạng phản tam giác:
\[
  m(i,j)+m(j,j')\le m(j,j)+m(i,j')=m(i,j').
\]
Gọi $k$ là một vị trí đạt cực tiểu trong công thức của $m(i,j')$, chọn vị trí
lớn nhất nếu có nhiều vị trí cùng đạt cực tiểu:
\[
  m(i,j')=m(i,k-1)+m(k,j')+w(i,j').
\]
Do tính đối xứng, không mất tính tổng quát, giả sử $k\le j$. Khi đó
\[
\begin{aligned}
  m(i,j)+m(j,j')
  &\le w(i,j)+m(i,k-1)+m(k,j)+m(j,j') \\
  &\le w(i,j')+m(i,k-1)+m(k,j)+m(j,j') \\
  &\le w(i,j')+m(i,k-1)+m(k,j') \\
  &=m(i,j').
\end{aligned}
\]
Dòng thứ hai dùng tính đơn điệu của $w$ theo đoạn; dòng thứ ba dùng giả thiết
quy nạp cho đoạn ngắn hơn.

\paragraph{Trường hợp 2: $i<i'<j<j'$.}
Gọi $y$ và $z$ lần lượt là các quyết định đạt cực tiểu lớn nhất của
$m(i',j)$ và $m(i,j')$:
\[
\begin{cases}
  m(i',j)=w(i',j)+m(i',y-1)+m(y,j),\\
  m(i,j')=w(i,j')+m(i,z-1)+m(z,j').
\end{cases}
\]
Do tính đối xứng, không mất tính tổng quát, giả sử $z\le y$. Theo định nghĩa,
ta có $i<z\le y\le j$. Vì vậy
\[
\begin{aligned}
  &m(i,j)+m(i',j')\\
  &\le w(i,j)+m(i,z-1)+m(z,j)
       +w(i',j')+m(i',y-1)+m(y,j') \\
  &\le w(i,j')+w(i',j)
       +m(i',y-1)+m(i,z-1)+m(z,j)+m(y,j') \\
  &\le w(i,j')+w(i',j)
       +m(i',y-1)+m(i,z-1)+m(y,j)+m(z,j') \\
  &=m(i,j')+m(i',j).
\end{aligned}
\]
Dòng thứ hai dùng bất đẳng thức tứ giác của $w$; dòng thứ ba dùng giả thiết
quy nạp cho $m$ trên các đoạn ngắn hơn. Từ hai trường hợp trên, theo quy nạp
toán học, hàm $m(i,j)$ cũng thỏa mãn bất đẳng thức tứ giác.

Ta định nghĩa $s(i,j)$ là giá trị lớn nhất của biến quyết định làm cho
$m(i,j)$ đạt cực tiểu:
\[
  s(i,j)=
  \max\left\{
    k\ \middle|\ 
    m(i,j)=w(i,j)+m(i,k-1)+m(k,j)
  \right\}.
  \tag{1.3}
\]

\paragraph{Định lý 2.}
Nếu $m(i,j)$ thỏa mãn bất đẳng thức tứ giác, thì $s(i,j)$ có tính đơn điệu:
\[
  s(i,j)\le s(i,j+1)\le s(i+1,j+1).
\]

\paragraph{Chứng minh.}
Do tính đối xứng, chỉ cần chứng minh
$s(i,j)\le s(i,j+1)$. Khi $i>j$, ta có
$s(i,j)=s(i,j+1)=\infty$, mệnh đề hiển nhiên đúng. Khi $i=j$, ta có
$s(i,j)=0\le s(i,j+1)$. Dưới đây xét trường hợp $i<j$.

Để tiện trình bày, ký hiệu $m_k(i,j)$ là giá trị hàm mục tiêu khi biến quyết
định lấy giá trị $k$:
\[
  m_k(i,j)=w(i,j)+m(i,k-1)+m(k,j).
\]
Khi đó $m_{s(i,j)}(i,j)=m(i,j)$.

Vì $m(i,j)$ thỏa mãn bất đẳng thức tứ giác, với mọi $k\le k'\le j$ ta có
\[
  m(k,j)+m(k',j+1)\le m(k',j)+m(k,j+1).
\]
Cộng vào hai vế lượng
\[
  w(i,j)+m(i,k-1)+w(i,j+1)+m(i,k'-1),
\]
ta được
\[
  m_k(i,j)+m_{k'}(i,j+1)
  \le m_k(i,j+1)+m_{k'}(i,j),
\]
tức là
\[
  m_k(i,j)-m_{k'}(i,j)
  \le
  m_k(i,j+1)-m_{k'}(i,j+1).
  \tag{1.4}
\]
Từ (1.4), suy ra
\[
  m_{k'}(i,j)\le m_k(i,j)
  \Longrightarrow
  m_{k'}(i,j+1)\le m_k(i,j+1).
  \tag{1.5}
\]

Vì $k'\ge k$, và với mọi $k<s(i,j)$ đều có
\[
  m_{s(i,j)}(i,j)=m(i,j)\le m_k(i,j),
\]
nên theo hệ quả (1.5), ta có
\[
  m_{s(i,j)}(i,j+1)\le m_k(i,j+1).
\]
Do đó quyết định tối ưu lớn nhất của $m(i,j+1)$ không thể nhỏ hơn $s(i,j)$,
tức là $s(i,j+1)\ge s(i,j)$. Định lý được chứng minh.

Sau khi biết biến quyết định $s(i,j)$ có tính đơn điệu, ta nhận được
\[
  s(i,j-1)\le s(i,j)\le s(i+1,j).
\]
Vì vậy công thức chuyển trạng thái (1.1) tương đương với
\[
  m(i,j)=
  \begin{cases}
    \displaystyle
    \min_{s(i,j-1)\le k\le s(i+1,j)}
      \{m(i,k-1)+m(k,j)+w(i,j)\}, & i<j,\\[0.8em]
    0, & i=j,\\
    \infty, & i>j.
  \end{cases}
  \tag{1.6}
\]
Giai đoạn của quy hoạch động là $l=j-i$. Khi cần tính $m(i,j)$, hai giá trị
$s(i,j-1)$ và $s(i+1,j)$ chắc chắn đã được tính trước. Do đó ta có thể thu
hẹp phạm vi liệt kê của biến quyết định $k$, từ đó giảm số phép tính.

Công thức (1.1) có độ phức tạp thời gian $O(n^3)$. Với công thức (1.6), độ
phức tạp là
\[
\begin{aligned}
  &O\left(
    \sum_{l=2}^{n}\sum_{i=1}^{n+1-l}
    \bigl(1+s(i+1,i+l-1)-s(i,i+l-2)\bigr)
  \right)\\
  &=
  O\left(
    \sum_{l=2}^{n}
    \bigl(n+1-l+s(n+2-l,n)-s(1,l-1)\bigr)
  \right)\\
  &\le
  O\left(
    \sum_{l=2}^{n} (2n+1-2l)
  \right)
  =
  O\bigl((n-1)^2\bigr).
\end{aligned}
\]
Vì vậy thuật toán sau khi tối ưu có độ phức tạp $O(n^2)$.

\section{Ứng dụng của bất đẳng thức tứ giác}

Thực ra, bất đẳng thức tứ giác không chỉ giới hạn ở dạng (1.1). Nhiều phương
trình chuyển trạng thái tương tự cũng thỏa mãn bất đẳng thức tứ giác.

\subsection*{Ví dụ: cây tìm kiếm nhị phân tối ưu}

\subsubsection*{Mô tả bài toán}

Từ $n$ phần tử
\[
  e_1<e_2<\cdots<e_n
\]
có thể tạo ra
\[
  \frac{1}{n+1}\binom{2n}{n}
\]
cây tìm kiếm nhị phân khác nhau. Với một cây tìm kiếm nhị phân $T$ cho trước,
giả sử nút chứa giá trị $e_i$ là $t_i$. Ta định nghĩa chi phí truy cập nút đó,
ký hiệu $c_i$, là số cạnh trên đường đi duy nhất nối gốc với nút đó. Đặc biệt,
chi phí truy cập nút gốc bằng $0$.

Cho thêm $n$ hằng số $f_1,f_2,\ldots,f_n$. Tổng trọng số của cây tìm kiếm nhị
phân $T$ được định nghĩa là
\[
  \operatorname{sum}_T=\sum_{i=1}^{n} f_i\cdot c_i.
\]
Ta gọi cây tìm kiếm nhị phân có tổng trọng số nhỏ nhất là cây tìm kiếm nhị
phân tối ưu. Cho $n$ và $f_1,f_2,\ldots,f_n$, hãy tìm tổng trọng số của cây
tối ưu.

\subsubsection*{Phân tích bài toán}

Điều kiện $e_1<e_2<\cdots<e_n$ xác định rõ tính có thứ tự của bài toán. Định
nghĩa
\[
  w(i,j)=
  \begin{cases}
    \displaystyle \sum_{k=i}^{j} f_k, & i\le j,\\[0.6em]
    0, & i>j.
  \end{cases}
\]
Ta nhận được phương trình chuyển trạng thái
\[
  m(i,j)=
  \begin{cases}
    \displaystyle
    \min_{i\le k\le j}
      \{m(i,k-1)+w(i,k-1)+m(k+1,j)+w(k+1,j)\}, & i<j,\\[0.8em]
    0, & i\ge j.
  \end{cases}
  \tag{2.1}
\]
Khi đó $m(1,n)$ chính là đáp án cần tìm.

\subsubsection*{Tối ưu hóa}

Rõ ràng công thức (2.1) có độ phức tạp thời gian $O(n^3)$. Tuy nhiên, chú ý
rằng (2.1) có dạng khá giống (1.1), nên ta thử chứng minh nó thỏa mãn bất
đẳng thức tứ giác. Từ đó có thể suy ra tính đơn điệu của biến quyết định và
giảm độ phức tạp xuống $O(n^2)$.

Trước hết, không khó thấy $w$ thỏa mãn tính đơn điệu theo quan hệ bao hàm của
đoạn và bất đẳng thức tứ giác. Dưới đây ta chứng minh $m$ cũng thỏa mãn bất
đẳng thức tứ giác. Để tiết kiệm dung lượng, một số chi tiết biên được lược
bớt.

Khi $i=i'$ hoặc $j=j'$, bất đẳng thức hiển nhiên đúng. Ta chứng minh quy nạp
theo $l=j'-i$.

\paragraph{Trường hợp 1: $i<i'=j<j'$.}
Bất đẳng thức trở thành
\[
  m(i,j)+m(j,j')\le m(i,j').
\]
Gọi $k$ là một quyết định đạt cực tiểu của $m(i,j')$:
\[
  m(i,j')=
  m(i,k-1)+w(i,k-1)+m(k+1,j')+w(k+1,j').
\]
Do tính đối xứng, không mất tính tổng quát, giả sử $k\le j$. Khi đó
\[
\begin{aligned}
  m(i,j)+m(j,j')
  &\le m(i,k-1)+w(i,k-1)+m(k+1,j)+w(k+1,j)+m(j,j')\\
  &\le m(i,k-1)+w(i,k-1)+m(k+1,j')+w(k+1,j')\\
  &=m(i,j').
\end{aligned}
\]

\paragraph{Trường hợp 2: $i<i'<j<j'$.}
Gọi $y$ và $z$ lần lượt là các quyết định đạt cực tiểu của $m(i',j)$ và
$m(i,j')$:
\[
\begin{cases}
  m(i',j)=
  m(i',y-1)+w(i',y-1)+m(y+1,j)+w(y+1,j),\\
  m(i,j')=
  m(i,z-1)+w(i,z-1)+m(z+1,j')+w(z+1,j').
\end{cases}
\]
Do tính đối xứng, không mất tính tổng quát, giả sử $z\le y$. Khi đó
$i\le z\le y\le j$, và
\[
\begin{aligned}
  &m(i,j)+m(i',j')\\
  &\le m(i,z-1)+w(i,z-1)+m(z+1,j)+w(z+1,j)\\
  &\quad +m(i',y-1)+w(i',y-1)+m(y+1,j')+w(y+1,j')\\
  &\le m(i,z-1)+w(i,z-1)+m(z+1,j')+w(z+1,j')\\
  &\quad +m(i',y-1)+w(i',y-1)+m(y+1,j)+w(y+1,j)\\
  &=m(i,j')+m(i',j).
\end{aligned}
\]
Tổng hợp lại, theo quy nạp toán học, mệnh đề được chứng minh.

Vì $m$ thỏa mãn bất đẳng thức tứ giác, ta có thể chứng minh biến quyết định
$s(i,j)$ có tính đơn điệu. Quá trình chứng minh giống với phần dành cho công
thức (1.1), nên không nhắc lại. Do đó, khi $i<j$, công thức (2.1) tương đương
với
\[
  m(i,j)=
  \min_{s(i,j-1)\le k\le s(i+1,j)}
  \{m(i,k-1)+w(i,k-1)+m(k+1,j)+w(k+1,j)\}.
\]
Phương trình chuyển trạng thái này chỉ cần thời gian $O(n^2)$.

\end{document}
